\documentclass[11pt]{article}
\usepackage[a4paper,margin=21mm]{geometry}
\usepackage{fontspec}
\setmainfont{Latin Modern Roman}
\newfontfamily\CJKfont{Noto Serif CJK SC}
\XeTeXlinebreaklocale "zh"
\XeTeXlinebreakskip=0pt plus 1pt
\usepackage{amsmath,amssymb,amsthm,amscd,graphicx,color}
\usepackage{xurl}
\usepackage[unicode,hidelinks]{hyperref}
\allowdisplaybreaks
\setlength\emergencystretch{3em}
\numberwithin{equation}{section}

%\usepackage{amsfonts,a4wide}
%\usepackage{enumerate}
%\usepackage[latin1]{inputenc}
\usepackage[all]{xy}
\usepackage{tikz-cd}
%\usepackage{cleveref}

\newcommand{\Sy}{\mathcal S}
\newcommand{\hh}{\mathbb{H}}
\newcommand{\D}{\mathcal{D}}
\newcommand{\J}{\mathcal{J}}
\newcommand{\Ab}{\mathcal{A}b}
\newcommand{\bbn}{\mathbb{N}}
\newcommand{\bbq}{\mathbb{Q}}
\newcommand{\R}{\mathbb{R}}
\newcommand{\RP}{\mathbb{RP}}
\newcommand{\CP}{\mathbb{CP}}
\newcommand{\HP}{\mathbb{HP}}
\newcommand{\GL}{\operatorname{GL}}
\newcommand{\Z}{\mathbb{Z}}
\newcommand{\ZZ}{\mathbb{Z}}
\newcommand{\Q}{\mathbb{Q}}
\newcommand{\QQ}{\mathbb{Q}}
\newcommand{\JJ}{\mathcal J}
\newcommand{\C}{\mathbb{C}}
\newcommand{\OO}{\mathbb{O}}
\newcommand{\HH}{\mathbb{H}}
\newcommand{\U}{\mathcal{U}}
\newcommand{\V}{\mathcal{V}}
\newcommand{\W}{\mathcal{W}}
\newcommand{\F}{\mathbb{F}}
\newcommand{\mbn}{\mathbf{n}}
\newcommand{\mbp}{\mathbf{p}}
\newcommand{\mbq}{\mathbf{q}}
\newcommand{\p}{\mathfrak p}
\newcommand{\q}{\mathfrak q}
\newcommand{\im}{\operatorname{Im}}
\newcommand{\mdc}{\operatorname{mdc}}
\DeclareMathOperator{\tr}{tr}
\DeclareMathOperator{\sgn}{sgn}
\DeclareMathOperator{\inter}{int}
\DeclareMathOperator{\Diff}{Diff}
\DeclareMathOperator{\ClImm}{ClImm}
\DeclareMathOperator{\Fun}{Fun}
\DeclareMathOperator{\coker}{coker}
\DeclareMathOperator{\PSh}{Psh}
\DeclareMathOperator{\Sp}{Sp}
\DeclareMathOperator{\Mod}{Mod}
\DeclareMathOperator{\Et}{\acute{E}t}
\DeclareMathOperator{\colim}{colim}
\DeclareMathOperator{\Hom}{Hom}
\DeclareMathOperator{\Ext}{Ext}
\DeclareMathOperator{\Ch}{Ch}
\DeclareMathOperator{\YExt}{YExt}
\DeclareMathOperator{\ev}{ev}
\DeclareMathOperator{\Supp}{Supp}
\DeclareMathOperator{\Tot}{Tot}
\DeclareMathOperator{\Spec}{Spec}
\DeclareMathOperator{\Tor}{Tor}
\DeclareMathOperator{\Der}{Der}
\DeclareMathOperator{\id}{id}
\DeclareMathOperator{\Gr}{Gr}
\DeclareMathOperator{\Op}{Op}
\DeclareMathOperator{\Emb}{Emb}
\DeclareMathOperator{\SympEmb}{SympEmb}
\DeclareMathOperator{\SympImm}{SympImm}
\DeclareMathOperator{\Th}{Th}
\DeclareMathOperator{\Map}{Map}
\DeclareMathOperator{\Aut}{Aut}
\DeclareMathOperator{\Id}{Id}
\DeclareMathOperator{\re}{Re}
\DeclareMathOperator{\Sq}{Sq}
\DeclareMathOperator{\End}{End}
\DeclareMathOperator{\Spin}{Spin}
\newcommand{\join}{\ast_{\Z/p}}
\newcommand{\cl}{\mathcal{C}{\rm l}}
\newcommand{\Cl}{\mathcal{C}{\rm l}}
\newcommand{\CCl}{\mathbb{C}{\rm l}}

\numberwithin{equation}{section}

\newtheorem{Theorem}{Theorem}[section]
\newtheorem*{Theorem*}{Theorem}
\newtheorem{Corollary}[Theorem]{Corollary}
\newtheorem{Lemma}[Theorem]{Lemma}
\newtheorem{Proposition}[Theorem]{Proposition}
 { \theoremstyle{definition}
\newtheorem{Definition}[Theorem]{Definition}
\newtheorem{Note}[Theorem]{Note}
\newtheorem{Example}[Theorem]{Example}
\newtheorem{Remark}[Theorem]{Remark} }

\begin{document}
\begin{center}\Large Rational minimal models of almost-complex-structure components on closed connected six-manifolds with vanishing first Betti number and nonzero c1c2-minus-c3 number\end{center}
\noindent\textbf{Stable ID:} 1307\quad\textbf{Author:} Gustavo Granja; Aleksandar Milivojević\par
\noindent\textbf{Work:} Topology of Almost Complex Structures on Six-Manifolds\par
\noindent\textbf{Source:} \url{https://sigma-journal.com/2022/093/}\par
\noindent\textbf{Version:} Official SIGMA 18 (2022) article 093, DOI 10.3842/SIGMA.2022.093; journal PDF and native TeX acquired 2026-09-30, exact bytes pinned by SHA256\par
\noindent\textbf{Rights:} CC BY 4.0. Authors retain copyright. \url{https://creativecommons.org/licenses/by/4.0/}. Reuse and typesetting adaptation; original language and mathematical content preserved.\par
\noindent\textbf{Difficulty (prior selection preserved):} {\CJKfont H2: The proof uses the canonical positive spinor bundle to replace the almost-complex section component by a circle quotient of sphere-bundle sections. It computes the rationalized section space, uses the nonzero Chern number to force the degree-one transgression, and eliminates the resulting acyclic generator pair from a Sullivan model. This ties spinorial geometry to the fundamental group and the full rational homotopy model of a function-space component; the base manifold itself need not be nilpotent.}\par
\medskip\noindent\textbf{Editorial boundary note (not author text):} Selected complete Theorem4.4 for the component of a chosen almost complex structure J on closed connected six-manifold M with b1(M)=0 and integral(c1c2-c3)!=0. No nilpotency assumption on M, integrability of J or all-component connectedness is made. Full singular-cohomology/graded-algebra/section-space conventions, spinorial/Clifford local supports2.2–2.3 and3.1–3.3, canonical section/lift, full S7-section/S1-quotient fibration, actual finite cyclic fundamental-group calculation and explicit model generator elimination are retained. Prior section-space rationalization, Thom mapping-space calculation and Crabb–Sutherland fundamental-group theorem stay exact established references. Intersection and four-dimensional results remain context and receive no counts. bbm/shadow imports have zero literal callers in the retained body and are omitted without notation replacement. Native commutative diagrams and all signs/orientations/generator multiplicities remain editable and unchanged.\par
\medskip\hrule\medskip\noindent\textbf{Author text begins below.}\par

% Original source lines 141--144
Let $M$ denote an oriented six-manifold and $\mathcal J(M)$ be the space of
all almost complex structures on $M$ inducing the given orientation.
Fixing a Riemannian metric on $M$, the inclusion $J(M) \hookrightarrow \mathcal J(M)$ of the subspace of almost complex structures which are orthogonal with respect to the metric is a homotopy equivalence, allowing us to analyze $\mathcal J(M)$ via a space of sections of a fiber bundle with compact fibers. Generally, over an oriented Riemannian $2n$-manifold $M$ one can consider the ${\rm SO}(2n)/{\rm U}(n)$ bundle
$Z_+(M) \to M$ whose fiber over any $x \in M$ is the set of (linear) orthogonal complex structures on $T_x M$ compatible with the orientation on $M$; the space $Z_+(M)$ is known as the \emph{$($positive$)$ twistor space} of $M$. Sections of this bundle correspond to orthogonal almost complex structures on $M$ compatible with the orientation.


% Original source lines 156--172
Our second main result is a description of the rational homotopy type of (a given component of) the space of almost complex structures on six-manifolds, under the additional assumption that $b_1(M) = 0$ and $\int_M c_1c_2 - c_3 \neq 0$ (Theorem~\ref{Qhomotopy6}). For this we use an expression of $J(M)$ as the quotient of the space of sections of an $S^7$-bundle by an $S^1$-action (Theorem~\ref{fibration}), which one can compare with the result obtained for $c_1=0$ in~\cite[Proposition~5.1]{FGM21}.
Throughout we provide examples for all of the mentioned results.

\subsection*{Notation and conventions}

 The projectivization of a complex vector bundle $E$ is denoted by $\mathbb{P}(E)$. We denote the Chern classes of a complex vector bundle by~$c_i$, and the Pontryagin classes of a real vector bundle by~$p_i$; if the vector bundle is complex, by its Pontryagin classes we mean those of the underlying real bundle. The Euler class is denoted by~$\mathsf{e}$. The Euler characteristic of a space is denoted by~$\chi$, and~$\sigma$ is the signature of an oriented closed manifold. All manifolds are assumed connected unless otherwise stated.

We denote by $\Gamma(E)$ the space of sections of a bundle $E \to M$ (where $E$ is not necessarily a~vector bundle). If $s \colon M \to E$ is a section, $\Gamma(E)_s$ denotes the connected component of~$s$ in~$\Gamma(E)$.
For $E \to M$ a vector bundle with a metric, we write $S(E) \to M$ for the associated sphere bundle. We denote by $J(M)$ the connected component of a specific point in $\Gamma(Z_+(M))$, which will be clear from context. For a vector space~$V$ with an inner product, $J(V)$ will denote the space of all orthogonal (linear) complex structures on $V$. If $V$ is oriented, $J_+(V)$ (respectively~$J_-(V)$) denotes the connected component of~$J(V)$ consisting of elements which induce the given orientation on $V$ (respectively the opposite orientation). By $\Map(X,Y)$ we denote the space of unbased maps $X \to Y$ with the compact-open topology.

Cohomology is singular cohomology with integral coefficients unless another choice of coefficients is indicated. The evaluation of a cohomology class $\alpha$ on a homology class $A$ is denoted by $\langle \alpha, A \rangle$ or by $\int_A \alpha$. In the context of rational homotopy theoretic models, $\Lambda(x_i)$ denotes the free graded-commutative algebra on generators $x_i$.

We will make frequent use of the projective bundle formula, which states that for a complex vector bundle $E \to M$ of complex rank $k$, the cohomology of the projectivized bundle $\mathbb{P}(E) \xrightarrow{p} M$ is the free $H^*(M)$-algebra on one generator $x$ subject only to the relation $x^k + c_1(E)x^{k-1} + \cdots + c_{k-1}(E)x + c_k(E) = 0$, where $x$ is the first Chern class of the line bundle $\mathcal{O}_{E^*}(1)$ (see~\cite[Definition~15.13]{D12}). This line bundle is dual
to the tautological (Hopf) line subbundle of $p^*(E)$ whose fiber over $\ell \in \mathbb{P}(E_x)$ is $\ell \subset E_x$. It restricts to $\mathcal{O}(1)$ on each fiber $\CP^{k-1}$.

For a real vector space $V$ with an endomorphism $J$ squaring to $-\mathrm{Id}$, we denote by $\Lambda_\C^*(V)$ the (complex) exterior algebra of the complex vector space $(V,J)$. For an almost complex manifold $(M,J)$, $\Lambda_\C^\ast (TM)$ is the (complex) exterior bundle associated to $(TM,J)$.


\setcounter{section}{1}
% Original source lines 173--335
\section{Preliminaries}\label{section2}
In this section we review some of the models for the space of orthogonal complex structures on a Euclidean space following~\cite[Section~IV.9]{LM}, to which we refer for further details.

Let $(V, (\cdot, \cdot ))$ be a $2n$-dimensional real inner product space. We will write \[J(V)=\big\{ J \in \End(V) \colon J^2=-I, \, J \text{ is orthogonal} \big\}\]
for the space of orthogonal complex structures on $V$. This space naturally has the structure of a complex algebraic variety, given by the isomorphism
\[ J(V) \xrightarrow{\phi} \Gr_n^{\rm Iso}(V\otimes \C) \]
with the Grassmannian of $n$-dimensional complex subspaces of $V\otimes \C$, which are isotropic with respect to the complex bilinear form on $V\otimes \C$ obtained from the inner product on $V$ by linear extension. The map $\phi$ is defined by
\[ \phi(J) = \{ v\otimes 1 + Jv \otimes {\rm i} \colon v \in V\} \]
(it assigns to $J$ the plane $V_J^{0,1} \subset V\otimes \C$, which is the $-{\rm i}$ eigenspace of the complex linear extension of $J$ to $V\otimes \C$).

There is a tautologous complex vector bundle on $J(V)$ defined by
\[J(V) \times V \xrightarrow{\pi_1} J(V),\]
where the fiber over $J\in J(V)$ is given the complex structure determined by $J$.
The isomorphism $\phi$ is covered by a bundle map to the dual (or conjugate) tautological bundle over the Grassmannian.

\begin{Lemma}\label{positchern}
Consider the universal bundle
\[
J\big(\R^{2n}\big)={\rm O}(2n)/{\rm U}(n) \xrightarrow{\iota} {\mathcal B}{\rm U}(n) \to {\mathcal B}{\rm O}(2n).
\]
Then
\begin{enumerate}\itemsep=0pt
\item[$(i)$] $\iota$ classifies the tautologous complex vector bundle on $J\big(\R^{2n}\big)$.
\item[$(ii)$] $\iota^*(c_1)$ generates a subgroup of index $2$ in the infinite cyclic second cohomology group of
each of the two components of $J\big(\R^{2n}\big)$.
\item[$(iii)$] $\iota^*(c_1)$ is positive. In particular, $\iota^*(c_1)^{\frac{n^2-n}{2}}$ evaluates on the orientation class of each of the two components of $J\big(\R^{2n}\big)$ to a positive integer.
\end{enumerate}
\end{Lemma}
\begin{proof}
$(i)$ The universal bundle can be realized by taking ${\mathcal B}{\rm U}(n)={\mathcal E}{\rm O}(2n)/{\rm U}(n)$. The universal bundle over ${\mathcal B}{\rm U}(n)$ is then given by ${\mathcal E}{\rm O}(2n)\times_{{\rm U}(n)} \R^{2n} \to {\mathcal E}{\rm O}(2n)/{\rm U}(n)$, and pulls back along the inclusion of the fiber $\iota$ to ${\rm O}(2n) \times_{{\rm U}(n)} \R^{2n} \to {\rm O}(2n)/{\rm U}(n)$, which is the homogeneous expression of the tautologous bundle.

$(ii)$ The inclusion-induced map ${\mathcal B}{\rm U}(n) \to {\mathcal B}{\rm O}(2n)$ factors through the
double cover ${\mathcal B}{\rm SO}(2n) \allowbreak \to {\mathcal B}{\rm O}(2n)$, and the two components of ${\rm O}(2n)/{\rm U}(n)$ are,
respectively, the fiber ${\rm SO}(2n)/{\rm U}(n)$ of ${\mathcal B}{\rm U}(n) \to {\mathcal B}{\rm SO}(2n)$ and the image
of this fiber under a lift to ${\mathcal B}{\rm U}(n)$ of the nontrivial deck transformation of
${\mathcal B}{\rm SO}(2n)$. Therefore it suffices to prove the statement for the component
${\rm SO}(2n)/{\rm U}(n)$. This follows from the Serre spectral sequence of the bundle
${\mathcal B}{\rm U}(n) \to {\mathcal B}{\rm SO}(2n)$ (see for instance~\cite[Theorem~III.6.11]{MiTo}).

$(iii)$ $\iota^*(c_1)$ is the pullback of $c_1$ of the dual tautological bundle under the embedding
\[J\big(\R^{2n}\big) \xrightarrow{\cong} \Gr_n^{\rm Iso}\big(\C^{2n}\big) \subset \Gr_n\big(\C^{2n}\big).\]
The $n^{\mathrm{th}}$ exterior power of the dual tautological bundle on $\Gr\big(\C^{2n}\big)$ gives the
Pl\"ucker embedding of the Grassmannian in projective space. Hence $\iota^*(c_1)$ is the first
Chern class of a line bundle on $J\big(\R^{2n}\big)$ whose sections embed $J\big(\R^{2n}\big)$ in projective space and is therefore positive.\qedhere
\end{proof}

The above identification of orthogonal complex structures with isotropic planes in
$V\otimes \C$ leads to another description of the space $J(V)$ which will play an
important role in this paper.


Let $\Cl(V)$ denote the Clifford algebra determined by the quadratic form $q(v)=-\|v\|^2$, and let $\CCl(V)=\Cl(V)\otimes_\R \C$
denote its complexification. Let $S_{\C}$ denote an irreducible ($2^n$-dimensional) $\CCl(V)$-module.

For each element, i.e., \emph{spinor}, $\sigma \in S_{\C}$ we have the map $V\otimes_\R \C \to S_{\C}$ given by right Clifford multiplication by $\sigma$. A spinor is said to be \emph{pure} if the associated kernel is half-dimensional. Let ${\mathcal P}S _{\C} \subset S_\C$ denote the subset of pure spinors. As the subspace $\ker (\cdot \sigma) \subset V\otimes \C$ is isotropic
there is a natural map from the projectivization of the set of pure spinors to the isotropic Grassmannian
\[
\mathbb{P}( {\mathcal P}S _\C ) \xrightarrow{\psi} \Gr^{\rm Iso}_{2n}(V\otimes \C)
\]
defined by $[\sigma] \mapsto \ker (\cdot \sigma)$. This is a map of algebraic varieties and is in fact an ${\rm SO}(2n)$-equivariant isomorphism~\cite[Proposition~IV.9.7]{LM}.

The space $J(V)$ has two connected components corresponding to the two possible orientations induced by the complex structure. A choice of orientation for $V$ leads to a decomposition
\[J(V) = J_+(V) \coprod J_-(V)\]
with $J_+(V)$ the component corresponding to the given orientation. In turn this decomposes the Grassmannian of isotropic subspaces into the components of positive and negative isotropic subspaces.

On the level of spinors this decomposition takes the following form.
The complex volume element in $\CCl(V)$ is the element
$\omega_\C = {\rm i}^n e_1 \cdots e_{2n}$, where
$\{e_i\}$ is any oriented orthonormal basis. As $\omega_{\C}^2=1$, the volume element decomposes
the module $S_C$ into a direct sum of $+1$ and $-1$ eigenspaces
\[
S_\C = S_\C^+ \oplus S_\C^-.
\]
The elements of the summands are called the positive and negative spinors, respectively.
Note that, since $\omega_\C$ anti-commutes with the action of $V \subset \CCl(V)$, both
$S_\C^+$ and $S_\C^-$ are ${\rm spin}^{\rm c}$ representations, called the positive and negative spinor representations, respectively.

Every pure spinor is necessarily either positive or negative and the isomorphisms $\phi$, $\psi$ give rise to isomorphisms of algebraic varieties
\[
\mathbb{P}\big({\mathcal P}S _\C^\pm\big) \cong J_\pm(V).
\]
In dimensions $2n=4,6$ all nonzero spinors are pure~\cite[Remark~IV.9.12]{LM} and the isomorphisms identify the spaces
of orthogonal complex structure with the disjoint union of the projectivization of positive
and negative spinors.

A point $J \in J(V)$ gives rise to a useful description (cf.~\cite[Chapter I, equation~(5.27)]{LM})
of the irreducible module $S_\C$.
We will write $\langle \cdot, \cdot \rangle$ for the Hermitian inner product on the complex vector space $(V,J)$ whose real part is $ (\cdot, \cdot )$.
Let $\Lambda_\C^*(V)$ denote the exterior algebra of the complex vector space $(V,J)$.
Contraction of a vector $v \in V$ with $\omega \in \Lambda_\C^k(V)$ is determined by the expression
\[
v \lrcorner (w_1 \wedge \cdots \wedge w_k) = \sum_{j=1}^k (-1)^{j+1}\langle v, w_j \rangle w_1 \wedge \cdots \wedge \widehat{w_j} \wedge \cdots \wedge w_k.
\]
The action $V \otimes_\R \Lambda_\C^*(V) \to \Lambda_\C^*(V)$ defined by
\begin{equation*}
v \cdot \omega = v \wedge \omega - v \lrcorner \, \omega
\end{equation*}
satisfies
\begin{equation*}
v\cdot(v \cdot \omega) = - \|v\|^2 \omega,
\end{equation*}
as one easily checks using an orthonormal basis for $V$ having a
real multiple of a nonzero $v$ as its first element. The complex linear extension of this action gives $\Lambda_\C^*(V)$ the structure of a $\CCl(V)$-module of complex dimension $2^n$. This is the dimension of the irreducible complex $\CCl(V)$-module, so $\Lambda_\C^*(V)$ is a convenient description of this module.

\iffalse
Indeed, one can consider an orthonormal basis $v_1,Jv_1, \ldots, v_n, Jv_n$ for $V$
with $v=\alpha v_1$ and check that \eqref{clifalg} holds for each $\omega$ in the standard orthonormal basis for $\Lambda_\C^*(V)$:
$$
v \cdot v_{i_1} \wedge \cdots \wedge v_{i_k} = \begin{cases}
- \alpha \ v_{i_2} \wedge \cdots \wedge v_{i_k} & \text{ if } i_1=1\\
\alpha \ v_{1} \wedge v_{i_1} \wedge \cdots \wedge v_{i_k} & \text{ if } i_1 \neq 1
\end{cases}
$$
and hence
$$
v \cdot( v \cdot v_{i_1} \wedge \cdots \wedge v_{i_k} ) =
- \alpha^2 \ v_{i_1} \wedge v_{i_2} \wedge \cdots \wedge v_{i_k} = - \|v\|^2 v_{i_1} \wedge \cdots \wedge v_{i_k}
$$
\fi

The following is an easy computation which we include for convenience.

\begin{Lemma}\label{plusmin}
Let $(V,J)$ be an oriented $2n$-dimensional Euclidean real vector space with an orthogonal complex structure $J$ compatible with the orientation and let $S_\C=\Lambda^*_\C(V)$ be the irreducible $\CCl(V)$-module described above. Then
\[
S_\C^+ = \bigoplus_{k \text{ even }} \Lambda_\C^k(V), \qquad S_\C^- = \bigoplus_{k \text{ odd }} \Lambda_\C^k(V).
\]
\end{Lemma}
\begin{proof}
Let $e_1, \ldots, e_n$ be an orthonormal basis for $(V,J)$. Then we can write the complex volume element as
\[
\omega_\C = {\rm i}^n e_1 Je_1 \cdots e_n Je_n .
\]
Consider the basis $e_{i_1} \wedge \cdots \wedge e_{i_k}$ for $\Lambda^k_{\C}$ with $1\leq i_1 < \cdots < i_k \leq n$. Let $1 \leq m \leq n$
and assume that $m \not \in \{i_1, \ldots, i_k \}$. Then
\[
\begin{aligned}
(e_m Je_m) \cdot e_{i_1} \wedge \cdots \wedge e_{i_k} & = e_m \cdot \big( Je_m \wedge e_{i_1} \wedge \cdots \wedge e_{i_k} \big) \\
& = -\big\langle e_m, Je_m \big\rangle e_{i_1} \wedge \cdots \wedge e_{i_k} = -i \ e_{i_1} \wedge \cdots \wedge e_{i_k}.
\end{aligned}
\]
On the other hand, if $m \in \{i_1, \ldots, i_k\}$, let $l$ be such that $m=i_l$. Then
\[
\begin{aligned}
(e_m Je_m) \cdot e_{i_1} \wedge \cdots \wedge e_{i_k} & = e_m \cdot \big( {-} (-1)^{l-1} \big\langle Je_m, e_{i_l} \big\rangle e_{i_1} \wedge \cdots \wedge \widehat{e_{i_l}} \wedge \cdots \wedge e_{i_k} \big)\\
& = e_m \wedge\big( (-1)^l (-{\rm i}) e_{i_1} \wedge \cdots \wedge \widehat{e_{i_l}} \wedge \cdots \wedge e_{i_k} \big) \\
& = (-1)^{l-1} (-1)^l(-{\rm i}) e_{i_1} \wedge \cdots \wedge e_{i_k} = {\rm i} e_{i_1} \wedge \cdots \wedge e_{i_k}.
\end{aligned}
\]
Hence
\[
\omega_\C \cdot e_{i_1} \wedge \cdots \wedge e_{i_k} = {\rm i}^n {\rm i}^k (-{\rm i})^{n-k} e_{i_1} \wedge \cdots \wedge e_{i_k} = (-1)^{2n-k} e_{i_1} \wedge \cdots \wedge e_{i_k},
\]
which completes the proof.
\end{proof}

\begin{Remark}
\label{lift}
For $(V,J)$ an oriented $2n$-dimensional Euclidean real vector space with an orthogonal complex structure $J$ compatible with the orientation, the line in $S^+_\C$ corresponding to $J\in J_+(V)$ is $\Lambda^0_\C(V)\subset S_\C^+$.
Indeed,
\[(v \otimes 1 + w \otimes {\rm i}) \cdot 1 = v+Jw = 0 \ \Leftrightarrow \ w=Jv. \]
\end{Remark}


\setcounter{footnote}{1}
% Original source lines 336--419
\section{Twistor spaces of four- and six-manifolds}\label{section3}

In this section we describe the integral cohomology of the twistor space of (Riemannian) almost complex six-manifolds, along with the integral cohomology of the twistor space of arbitrary oriented (Riemannian) four-manifolds. Finally we compute the Chern classes of the Atiyah--Hitchin--Singer almost complex structure on the latter.

\subsection{Spinors on almost complex four- and six-manifolds}

Let $M$ be an oriented Riemannian manifold and let $Z(M)$ denote the space of orthogonal almost complex structures on $M$, called the \emph{twistor space}. The space
$Z(M)$ has two components corresponding to the structures which induce the given orientation on $M$ or the opposite orientation. Recall, we
denote these by $Z_{\pm}(M)$, and refer to them as the \emph{positive} and \emph{negative} twistor spaces.

The Atiyah--Hitchin--Singer almost complex structure on $Z(M)$ is defined as follows: the Levi-Civita connection induces a connection on the bundle $Z(M) \to M$ splitting $TZ(M)$ into a~direct sum of horizontal and vertical subspaces
\[
TZ(M) =T^{\rm h} Z(M) \oplus T^{\rm v}Z(M).
\]
The almost complex structure on $T^{\rm v} Z(M)$ is determined by the algebraic variety structure of the fibers $J(T_x M)$ of $Z(M) \to M$ explained in Section~\ref{section2}. At a point $J_x$ in the fiber
over $x\in M$ the complex structure on $T^{\rm h}_{J_x} Z(M) = T_x M$ is given\footnote{The Eells--Salamon almost complex structure is obtained replacing the vertical component of the Atiyah--Hitchin--Singer almost complex structure by its negative; see~\cite[Section~3]{Sa84} for details.} by $J_x$.



A ${\rm spin}^{\rm c}$ structure on $M$~\cite[Appendix~D]{LM} yields a complex spinor
bundle $S_\C(M)\to M$ together with an action
\[
\CCl(TM) \otimes S_\C(M) \to S_\C(M).
\]
The discussion in the previous section then provides a one-to-one correspondence between projectivized pure spinors on $M$ and orthogonal complex structures inducing a given orientation:
\[
Z_{\pm}(M) \cong \mathbb{P}\big( {\mathcal P}S _{\C}^{\pm}(M)\big).
\]
Now suppose $M$ is given an orthogonal almost complex structure $J$ compatible with the orientation. Then $J$ gives rise to a canonical ${\rm spin}^{\rm c}$ structure on $M$ via the group homomomorphism ${\rm U}(n) \xrightarrow{\kappa} \Spin^{\rm c}(2n)$, which is the unique
lift as a group homomorphism in the following diagram{\samepage
\[
\begin{tikzcd}
& \Spin^{\rm c}(2n) = \big(\Spin(2n) \times S^1\big)/\{\pm(1,1)\} \ar[d, "\pi \times \delta"] \\
{\rm U}(n) \ar[ur,"\kappa"] \ar[r,"\iota \times \det"] & {\rm SO}(2n) \times S^1.
\end{tikzcd}
\]
Here $\iota$ denotes the inclusion, $\pi\big(\big[h, {\rm e}^{{\rm i}\theta}\big]\big) = [h]$, and $\delta\big(\big[h,{\rm e}^{{\rm i}\theta}\big]\big)={\rm e}^{2{\rm i}\theta}$.}

An element $g\in {\rm U}(n)$ is diagonalized by some orthonormal basis $(e_1,Je_1, \ldots, e_n, Je_n)$. If
\[g=\mathrm{diag}\big({\rm e}^{{\rm i}\theta_1},\ldots, {\rm e}^{{\rm i}\theta_n}\big),\]
then
\[
\kappa(g)= \left[ \prod_{k=1}^n \left( \cos \frac{\theta_k} 2 + {\rm i} \sin \frac{\theta_k}{2} e_k J e_k\right), {\rm e}^{{\rm i} \frac{\sum \theta_k}{2}} \right]
\]
 (cf.~\cite[equation~(D.10)]{LM}). Using
the formulas for the action of $e_i Je_i$ in the proof of Lemma~\ref{plusmin}, we see that the spinor action of $\kappa(g)$
on $\Lambda^*_{\C}\big(\R^{2n}\big)$, where $\R^{2n}$ is equipped with the complex structure given by $\mathrm{diag}\left( \left(\begin{smallmatrix} 0 & -1 \\ 1 & 0 \end{smallmatrix}\right), \ldots, \left(\begin{smallmatrix} 0 & -1 \\ 1 & 0 \end{smallmatrix}\right) \right)$, agrees with the standard action of $g$ on this space. We obtain the following result:

\begin{Proposition}\label{identification}
Let $M$ be an oriented Riemannian manifold with an orthogonal almost complex structure $J$ compatible with the orientation. Then we have the following isomorphisms of smooth fiber bundles over $M$:
\begin{enumerate}\itemsep=0pt
\item[$1.$] If $\dim(M)=4$, then
\[
Z_+(M) \cong \mathbb{P}\big( \C \oplus \Lambda^2_\C(TM)\big), \qquad Z_{-}(M) \cong \mathbb{P}( TM).
\]
\item[$2.$] If $\dim(M)=6$, then
\[
Z_+(M) \cong \mathbb{P}\big( \C \oplus \Lambda^2_\C(TM)\big), \qquad Z_{-}(M) \cong \mathbb{P}\big( TM \oplus \Lambda_\C^3(TM)\big).
\]
\end{enumerate}
\end{Proposition}

\begin{Remark}
 The Atiyah--Hitchin--Singer and Eells--Salamon almost complex structures on (either) twistor space do not agree with the natural almost complex structure on the projectivization. Indeed, the projection map from the projectivization is complex linear whilst this is not the case for the twistor projection.
\end{Remark}

The projective bundle formula now gives the following formulas for the cohomology rings of the components of the twistor space in terms of the
Chern classes of $(TM, J)$.

\begin{Corollary}
\label{cohomtwis}
Let $M$ be an oriented Riemannian manifold with an orthogonal almost complex structure $J$ compatible with the orientation. We have the following isomorphisms as algebras over $H^*(M)$:
\begin{enumerate}\itemsep=0pt
\item[$1.$] If $\dim(M)=4$, then
\begin{align*}
&H^*(Z_+(M))\cong H^*(M)[x]/\big(x^2+c_1 x\big), \\
&H^*(Z_{-}(M))\cong H^*(M)[x]/\big(x^2 + c_1 x + c_2\big).
\end{align*}
\item[$2.$] If $\dim(M)=6$, then
\begin{align*}
&H^*(Z_+(M))\cong H^*(M)[x]/\big(x^4 +2c_1x^3 +\big(c_1^2+c_2\big) x^2 + (c_1c_2-c_3)x\big), \\
&H^*(Z_{-}(M))\cong H^*(M)[x]/\big(x^4 + 2c_1x^3+ \big(c_1^2+c_2\big)x^2 + (c_1c_2+c_3)x\big). \end{align*}
\end{enumerate}
\end{Corollary}

\setcounter{section}{3}
% Original source lines 779--825
\section[On the homotopy type of the space of almost complex structures on six-manifolds]{On the homotopy type of the space\\ of almost complex structures on six-manifolds}

In this section we study the homotopy type of the components of the space of almost complex structures on connected six-manifolds $M$, complementing our previous treatment~\cite{FGM21} of the case when $c_1=0$.

\begin{Theorem}\label{fibration}
Let $M$ be an oriented Riemannian manifold of dimension four or six, $J$ an orthogonal almost complex structure on $M$ compatible with the orientation, and $J(M)$ its component in the space of orthogonal almost complex structures. Let $S_{\C}^+(M)$ be the positive spinor bundle on $M$ determined by $J$. Then the map
\[
\Gamma\big( S\big(S_\C^+(M)\big)\big) \to J(M)
\]
is surjective, and for any lift $s$ of $J$,
\[
\Map\big(M,S^1\big) \to \Gamma\big( S\big(S_\C^+\big)\big)_s \to J(M)
\]
is a fiber sequence.
\end{Theorem}
\begin{proof}
Identifying $S_\C^+(M)$ with $\oplus_{k \text{ even}} \Lambda^k_\C(TM)$, Remark~\ref{lift} shows that
the constant section $s \colon M \to \Lambda^0_\C(TM) \subset S_\C^+(M)$ defined by
$s(x)=1$ lifts the section $J$. As $S\big(S_\C^+\big) \to Z_+(M)$ is a~fiberwise fibration over $M$, the map induced on sections is a~fibration. Since the fiber over $J$ is nonempty, any choice of lift of $J$ identifies the
fiber with $\Map\big(M,S^1\big)$.
\end{proof}

\begin{Remark} The above argument remains valid in all dimensions if one replaces $S\big(S_{\C}^+(M)\big)$ with $S\big(S_{\C}^+(M)\big) \cap {\mathcal P}S ^+_{\C}(M)$. \end{Remark}

Now let $(M,J)$ be a closed almost complex Riemannian six-manifold with $H^1(M;\ZZ) = 0$ (equivalently, $b_1 = 0$) satisfying $c_1c_2 - c_3 \neq 0$. By Theorem~\ref{fibration}, we have the fibration $\Map\big(M,S^1\big) \to \Gamma\big( S\big(S_\C^+(M)\big)\big)_s \to J(M)$. Since $H^1(M;\ZZ) = 0$, evaluation at a chosen basepoint gives a homotopy equivalence $\Map\big(M,S^1\big) \xrightarrow{{\rm ev}} S^1$, so we have a principal fiber sequence $S^1 \to \Gamma\big(S\big(S_\C^+(M)\big)\big)_s \to J(M)$.

We can thus consider instead the fiber sequence \begin{equation}\label{fibration2} \Gamma(S\big(S_{\C}^+\big))_s \to J(M) \to {\mathcal B}S ^1 \simeq \CP^\infty. \end{equation} Now, $S\big(S_{\C}^+\big)$ is an oriented fiber bundle over $M$ with fiber $S^7$, and hence it is classified by a map $M \to {\mathcal B}{\rm Aut}^+\big(S^7\big)$ to the classifying space of orientation-preserving homotopy automorphisms of~$S^7$. It is known that ${\mathcal B}{\rm Aut}^+\big(S^7\big)_{\QQ} \simeq K(\QQ, 8)$, where the subscript of $\QQ$ denotes (Sullivan) rationalization~\cite[Section~11]{S77}. By~\cite[Theorem~5.3]{M87}, the rationalization of $\Gamma\big(S\big(S_{\C}^+\big)\big)_s$ is homotopy equivalent to (the connected component corresponding to $s$ in) the space of sections of the fiberwise rationalized $S^7$ bundle over $M$. Since $H^8(M;\QQ) = 0$, this latter bundle is trivial, and hence $\big(\Gamma\big(S\big(S_{\C}^+\big)\big)_s\big)_{\QQ}$ is homotopy equivalent to a connected component of $\Map\big(M, S^7_{\QQ}\big)$. Denoting the Betti numbers of $M$ by $b_i$, by Thom's theorem on the space of maps into an Eilenberg--Maclane space, $\Map\big(M, S^7_{\QQ}\big)$ is in fact connected and has the homotopy type of $S^1_{\QQ} \times S^7_{\QQ} \times K(\QQ, 3)^{b_4} \times K(\QQ, 4)^{b_3} \times K(\QQ,5)^{b_2}$.

We can describe the fundamental group of $J(M)$ by using a theorem of Crabb and Sutherland:

\begin{Theorem}[{\cite[Theorem~2.12(i)]{CS84}}] Let $M$ be a closed connected $2n$-manifold and $\xi$ a complex $(n+1)$-plane bundle over $M$. Denote by $N\xi$ any component of the space of sections of $\mathbb{P}(\xi)$ whose elements lift to sections of $\xi$. Then $\pi_1(N\xi)$ is a central extension \[0 \to \ZZ/\big(\big\langle c_n(\xi), [M] \big\rangle \big) \to \pi_1(N\xi) \to H^1(M) \to 0.\]
\end{Theorem}

By Theorem~\ref{fibration}, we can apply this to $\xi = S_{\C}^+(M) \cong \C \oplus \Lambda_{\C}^2 TM$, which satisfies ${N \xi = J(M)}$. Therefore, the fundamental group of $J(M)$ is given by $\ZZ$ modulo $\int_M\! c_1c_2-c_3$ (where ${c_i \! =\! c_i(TM)}$), see Corollary~\ref{cohomtwis}. By~\cite{M87}, $J(M)$ is a nilpotent space as it is the space of sections of a fibration with nilpotent fiber \big(namely $\CP^3$\big) over a finite-dimensional base. We may thus rationalize the fibration (\ref{fibration2}) to obtain the fibration \[\big(\Gamma\big(S\big(S_{\C}^+\big)\big)_s\big)_{\Q} \to J(M)_{\Q} \to K(\Q, 2).\] Consider the degree one rational class corresponding to the factor $S^1_{\QQ}$ in the fiber. If $c_1c_2 - c_3 \neq 0$, then since $\pi_1(J(M)_\Q) = 0$, this class must hit (a nonzero multiple of) the degree two generator in $K(\Q,2)$ in the Serre spectral sequence. A model for the total space is given by the tensor product of the base and fiber's models with a perturbed differential on the fiber generators~\cite[Section~4]{S77}; the previous sentence thus tells us that a model for $J(M)$ is of the form
\[\big(\Lambda\big(x_2, z_1, z_7, z_3^{i}, z_4^{j}, z_5^{k}\big),\, {\rm d}z_1 = x_2, \, {\rm d}z_7 \in (x_2), \, {\rm d}z_3^i \in (x_2), \, {\rm d}z_4^j \in (x_2),\, {\rm d}z_5^k \in (x_2)\big),
\] where $(x_2)$ denotes the ideal generated by $x_2$, and $i$, $j$, $k$ range over sets of size $b_4$, $b_3$, $b_2$, respectively.

Notice that such a model is not minimal, and in fact a minimal model is obtained by quotienting out the differential ideal generated by~$z_1$. Indeed, by the argument in~\cite[Proposition~2]{VS76} we see that the map $(\Lambda, {\rm d}) \to \big(\Lambda/(z_1, {\rm d}z_1)\Lambda, \bar{{\rm d}}\big)$ induces an isomorphism on cohomology.
Here by $(\Lambda,{\rm d})$ we denote the differential graded algebra displayed above, and by $\bar{{\rm d}}$ the induced differential on the quotient. We see that
\[
(\Lambda/(z_1, x_2)\Lambda, \bar{{\rm d}}) \cong \big(\Lambda\big(z_7, z_3^{i}, z_4^{j}, z_5^{k}\big), \bar{{\rm d}}=0\big),
\] where the right-hand side is minimal. To summarize, we have the following:

\begin{Theorem}\label{Qhomotopy6} Let $M$ be a connected closed six-manifold with $b_1 = 0$ equipped with an almost complex structure $J$ such that $\int_M c_1c_2-c_3 \neq 0$. Then the space of almost complex structures on~$M$ in the component of $J$ is a nilpotent space with finite cyclic fundamental group and minimal model given by $\big(\Lambda\big(z_7, z_3^{i}, z_4^{j}, z_5^{k}\big), {\rm d} = 0\big)$; here $i$, $j$, $k$ range over sets of size $b_4(M)$, $b_3(M)$, $b_2(M)$ $(=b_4(M))$, respectively. In particular, this space is formal.
\end{Theorem}

We emphasize that the above makes no nilpotency assumption on $M$. Alternatively, a model for (a given component of) the space of almost complex structures can be obtained with the Haefliger--Sullivan model for the space of sections of a fibration~\cite[Section~11]{S77},~\cite{H82}, applied to $\CP^3 \to Z_+(M) \to M$. Our approach above lets one quickly describe the minimal model of~$J(M)$ using the geometry of the setup.
\begin{thebibliography}{99}
\bibitem[7]{CS84}
Crabb M.C., Sutherland W.A., Function spaces and {H}urwitz--{R}adon numbers,
 \href{https://doi.org/10.7146/math.scand.a-12068}{\textit{Math. Scand.}} \textbf{55} (1984), 67--90.

\bibitem[8]{D12}
Demailly J.-P., Complex analytic and differential geometry, available at
 \url{https://www-fourier.ujf-grenoble.fr/~demailly/manuscripts/agbook.pdf}.

\bibitem[13]{FGM21}
Ferlengez B., Granja G., Milivojevic A., On the topology of the space of almost
 complex structures on the six sphere, \textit{New York J.~Math.} \textbf{27}
 (2021), 1258--1273, \href{https://arxiv.org/abs/2108.00750}{arXiv:2108.00750}.

\bibitem[15]{H82}
Haefliger A., Rational homotopy of the space of sections of a nilpotent bundle,
 \href{https://doi.org/10.2307/1999931}{\textit{Trans. Amer. Math. Soc.}} \textbf{273} (1982), 609--620.

\bibitem[17]{LM}
Lawson H.B., Michelsohn M.-L., Spin geometry, \textit{Princeton Math. Ser.},
 Vol.~38, Princeton University Press, Princeton, NJ, 1989.

\bibitem[20]{MiTo}
Mimura M., Toda H., Topology of {L}ie groups.~{I},~{II}, \textit{Transl. Math.
 Monogr.}, Vol.~91, \href{https://doi.org/10.1090/mmono/091}{Amer. Math. Soc.}, Providence, RI, 1991.

\bibitem[21]{M87}
M{\o}ller J.M., Nilpotent spaces of sections, \href{https://doi.org/10.2307/2000693}{\textit{Trans. Amer. Math. Soc.}}
 \textbf{303} (1987), 733--741.

\bibitem[23]{Sa84}
Salamon S., Harmonic and holomorphic maps, in Geometry Seminar ``{L}uigi
 {B}ianchi''~{II}~-- 1984, \textit{Lecture Notes in Math.}, Vol.~1164,
 \href{https://doi.org/10.1007/BFb0081912}{Springer}, Berlin, 1985, 161--224.

\bibitem[25]{S77}
Sullivan D., Infinitesimal computations in topology, \href{https://doi.org/10.1007/BF02684341}{\textit{Inst. Hautes
 \'Etudes Sci. Publ. Math.}} \textbf{47} (1977), 269--331.

\bibitem[27]{VS76}
Vigu\'e-Poirrier M., Sullivan D., The homology theory of the closed geodesic
 problem, \href{https://doi.org/10.4310/jdg/1214433729}{\textit{J.~Differential Geometry}} \textbf{11} (1976), 633--644.

\end{thebibliography}
\end{document}
